\documentclass{report}
\usepackage{amsmath,amssymb}
\setlength{\parindent}{0mm}
\setlength{\parskip}{1em}
\begin{document}
\begin{center}
\rule{6in}{1pt} \
{\large Veronica Mejia Bustamante \\
{\bf Implementation of Iterative Solvers for the Digital Tomosynthesis Problem in GPUs}}

Dept of Math & CS \\ Emory University \\ 400 Dowman Dr  \\ W401 \\ Atlanta \\ GA 30322
\\
{\tt vmejia@emory.edu}\end{center}

Tomosynthesis imaging provides a viable alternative to computed
tomography (CT) and has obtained significant interest from the medical
community as a means for diagnostic radiology and radiation therapy. In
digital tomosynthesis imaging, multiple projections of an object are
obtained along a small range of different incident angles in order to
reconstruct a 3D representation of the object. In this paper we discuss
the implementation details of the polyenergetic digital breast
tomosynthesis reconstruction algorithm in a GPU using OpenCL. We describe
three different algorithm implementations: a serial implementation, a GPU
implementation threaded by functionality of the model, and a GPU fused
kernel implementation which is threaded to increase performance,
throughput, and GPU utilization in the application. We show that the
explicit kernel fusion achieves significant speed-up in the
reconstruction process of a clinical size patient data set, from running
over 100X faster than the version threaded by functionality to 200X
faster than the serial approach.


\end{document}
